\documentclass[11pt]{article}
\usepackage[utf8]{inputenc}
\usepackage[T1]{fontenc}
\usepackage{geometry}
\geometry{letterpaper, margin=1in}
\usepackage{amsmath,amssymb,amsfonts,amsthm}
\usepackage{booktabs}
\usepackage{graphicx}
\usepackage{hyperref}
\usepackage{url}
\usepackage{microtype}
\usepackage{xcolor}
\usepackage{subcaption}

\hypersetup{
    colorlinks=true,
    linkcolor=blue!70!black,
    citecolor=green!50!black,
    urlcolor=blue!80!black
}

\title{\textbf{StateShift: Tracking State-Dependent Reasoning Recovery\\Across Post-Training}}

\author{
  \textbf{Sham Satish Thakare} \\
  Independent AI Researcher \\
  \texttt{github.com/shamddd} \\
}
\date{October 2026}

\begin{document}
\maketitle

\begin{abstract}
Standard evaluations of reinforcement-learning post-training in large language models emphasize aggregate benchmark accuracy (e.g., Pass@1), potentially obscuring behavioral changes conditional on intermediate reasoning state. We introduce \textbf{StateShift}, a controlled diagnostic evaluation framework and benchmark protocol for measuring target-transition recovery under matched counterfactual reasoning states. In a prospective controlled experiment on mathematical reasoning ($N=454$ decontaminated competition problems, $29,056$ primary rollouts across the DeepScaleR-7B lineage), we isolate the \emph{State-Matched Interaction Contrast} ($\Gamma_t$), which differences out baseline problem difficulty and forward proficiency gains. Between the un-tuned base model ($t=0$) and the terminal reinforcement-learning checkpoint ($t=256$), we observe a statistically significant interaction gain of $\Gamma_{256} = \mathbf{+11.76\text{ percentage points}}$ ($\mu_{R,0}=38.34\% \to \mu_{R,256}=70.39\%$ vs.\ $\mu_{C,0}=38.92\% \to \mu_{C,256}=59.21\%$; $95\%$ problem-blocked bootstrap CI $[+9.55, +14.00]\text{ pp}$, cluster-robust $p < 10^{-24}$). On a strictly decontaminated sensitivity subset ($N_{\text{Strict}}=388$), the interaction remains stable at $\Gamma_{256,\text{Strict}} = \mathbf{+11.60\text{ pp}}$ ($p < 10^{-20}$), refuting benchmark leakage. Problem-level variance decomposition reveals that recovery gains are decoupled from clean accuracy gains ($r = 0.0794, p = 0.091, R^2 = 0.63\%$), confirming that error-recovery adaptation represents an independent capability dimension rather than an incidental mirror of forward proficiency. Two One-Sided Tests (TOST) confirm sequence-length equivalence within $\pm 1.0$ token ($p = 0.0256$), disproving test-time search expansion. Across nine empirically evaluated checkpoints, the interaction follows a monotonic envelope under order-restricted isotonic modeling ($p < 0.0001$) despite intermediate sampling variance; early emergence at Step 32 is nominally suggestive ($\Gamma_{32} = +3.33\text{ pp}, p_{\text{unadj}} = 0.019$), while family-wise significance is robustly confirmed by Step 96 ($\Gamma_{96} = +7.74\text{ pp}, p_{\text{Bonf}} < 0.001$). In an orthogonal ecological validation of $3,200$ unperturbed rollouts, models exhibit an endogenous error incidence of $18.19\%$ and autonomously self-correct in $30.93\%$ of error episodes. These results demonstrate that outcome-based reinforcement learning implicitly reshapes state-conditioned transition dynamics while leaving a persistent $29.61\%$ terminal failure regime.
\end{abstract}

\section{Introduction}
Reinforcement learning with verifiable rewards (RLVR) has emerged as a primary paradigm for scaling mathematical and algorithmic reasoning in foundation language models~\cite{Shao2024Deepseekmath, DeepSeek2025Deepseek, Luo2025Deepscaler}. Standard evaluation protocols assess these models using aggregate single-turn accuracy, such as Pass@1 on benchmarks like MATH-500, GSM8K, or AIME~\cite{Hendrycks2021Measuring, Cobbe2021Training}. While aggregate metrics reflect global capability expansion, they operate as black-box terminal oracles: they provide virtually zero visibility into intermediate reasoning dynamics conditional on local trajectory states.

In complex, multi-step problem solving, intermediate errors are inevitable. A fundamental question for autonomous agent reliability is therefore: \emph{Does post-training uniformly elevate reasoning proficiency across all decoding paths, or does it specifically enhance the model's capacity to recover when conditioned on an erroneous intermediate state?}

Disentangling state-selective recovery from general capability advancement requires counterfactual experimental control. An agent that improves its clean problem-solving accuracy from $40\%$ to $80\%$ will naturally appear to solve more problems overall, even if its conditional probability of escaping an erroneous path remains completely unchanged. Standard benchmark scores conflate these two phenomena.

\begin{figure}[t]
\centering
\includegraphics[width=0.88\linewidth]{figures/figure1_framework.pdf}
\caption{\textbf{The StateShift diagnostic evaluation framework.} For each problem instance, continuations are evaluated under matched Recovery ($R$, invalid intermediate deduction) and Control ($C$, length-matched valid deduction) conditions across post-training checkpoints to isolate the State-Matched Interaction Contrast ($\Gamma_t$).}
\label{fig:framework}
\end{figure}

To address this challenge, we introduce \textbf{StateShift}, a controlled experimental protocol and diagnostic benchmark suite designed to quantify state-dependent transition recovery. As shown in Figure~\ref{fig:framework}, StateShift pairs each problem with two counterfactual prefixes: an invalid intermediate prefix ($R$, Recovery condition) and a length-matched valid prefix ($C$, Control condition). By evaluating trajectories across reinforcement learning checkpoints under both conditions, StateShift isolates the \emph{State-Matched Interaction Contrast} ($\Gamma_t$).

Our primary empirical investigation tracks the open-source DeepScaleR RL post-training lineage~\cite{Luo2025Deepscaler} across nine checkpoints on $454$ decontaminated competition problems ($29,056$ primary rollouts). We observe a statistically significant interaction gain of $\Gamma_{256} = \mathbf{+11.76\text{ percentage points}}$ ($p < 10^{-24}$). Furthermore, we establish the ecological validity of this phenomenon in Study B by demonstrating that natural, unprompted decoding exhibits an $18.19\%$ Natural Error Incidence ($\text{NEI}$) and a $30.93\%$ Natural Recovery Rate ($\text{NRR}$).

\section{Related Work}

\paragraph{Self-Correction Capabilities in Language Models.}
Prior investigations into self-correction present conflicting conclusions. Huang et al.~\cite{Huang2023Can} and Stechly et al.~\cite{Stechly2024Large} demonstrated that frozen language models cannot intrinsically self-correct reasoning errors through multi-turn prompting without external oracle feedback. When prompted to review their solutions, models frequently alter correct derivations into incorrect ones. Conversely, SCoRe~\cite{Kumar2024Training} explicitly trains multi-turn policies using reinforcement learning designed to reward second-turn revisions. StateShift occupies a distinct paradigm: we evaluate single-rollout, continuous decoding in models trained \emph{only on final outcome rewards}, measuring emergent transition dynamics without multi-turn scaffolding or external verifier prompting.

\paragraph{Process Supervision and Intermediate States.}
Process Reward Models (PRMs) provide step-level supervision by scoring individual reasoning deductions~\cite{Lightman2023Let, Uesato2022Solving}. However, PRMs require training separate neural reward models that are prone to calibration drift and reward hacking. In contrast, StateShift utilizes deterministic symbolic equivalence oracles (SymPy) on verified mathematical domains, avoiding reward model noise.

\paragraph{Reinforcement Learning from Verifiable Rewards (RLVR).}
Recent post-training methodologies scale reasoning using Group Relative Policy Optimization (GRPO)~\cite{Shao2024Deepseekmath, DeepSeek2025Deepseek, Luo2025Deepscaler}. While these works document impressive jumps in benchmark accuracy, our study is the first to track the state-conditioned recovery trajectory across sequential checkpoints under strict test-time compute controls.

\section{The StateShift Evaluation Framework}

\subsection{Problem Formulation and Experimental Unit}
Let $\mathcal{P} = \{p_i\}_{i=1}^N$ denote a finite population of mathematical reasoning problems. For each problem $p_i$, let $Q_i$ represent the problem statement and $A_i^*$ the ground-truth final answer. An intermediate solution prefix $P_{i,s}$ represents a sequence of reasoning steps leading to state $s \in \{R, C\}$:
\begin{enumerate}
    \item \textbf{Recovery Prefix ($P_{i,R}$)}: An intermediate reasoning prefix terminating in an adjudicated, mathematically invalid deduction.
    \item \textbf{Control Prefix ($P_{i,C}$)}: A matched valid intermediate reasoning prefix matched for step depth and token length ($|\text{len}(P_{i,R}) - \text{len}(P_{i,C})| \le 3$ tokens).
\end{enumerate}

Given model checkpoint $\theta_t$ and prefix $P_{i,s}$, the model generates completion continuation $Y_{i,s,t,k}$ under stochastic sampling ($T=0.7$, Top-$p=0.95$). The binary outcome indicator is defined by deterministic symbolic verification:
\begin{equation}
y_{i,s,t,k} = \mathbb{I}\left(\text{Verify}\left(\text{Extract}(Y_{i,s,t,k}), A_i^*\right) = 1\right)
\end{equation}

\subsection{The State-Matched Interaction Contrast \texorpdfstring{($\Gamma_t$)}{(Gamma\_t)}}
Let $\mu_{s,t}$ denote the expected target-transition success probability for condition $s$ at checkpoint $t$:
\begin{equation}
\mu_{s,t} = \mathbb{E}_{i,k}[y_{i,s,t,k}] = \frac{1}{N K} \sum_{i=1}^N \sum_{k=1}^K y_{i,s,t,k}
\end{equation}
The marginal post-training gain under condition $s$ relative to baseline $\theta_0$ is $\Delta_{s,t} = \mu_{s,t} - \mu_{s,0}$. The State-Matched Interaction Contrast is formally defined as:
\begin{equation}
\Gamma_t \triangleq \Delta_{R,t} - \Delta_{C,t} = (\mu_{R,t} - \mu_{R,0}) - (\mu_{C,t} - \mu_{C,0})
\end{equation}

\paragraph{Two-Way Fixed-Effects Representation.}
We cast StateShift as a within-subject two-way fixed-effects linear probability model nested within problem blocks:
\begin{equation}
y_{i,s,t,k} = \alpha_i + \beta \cdot \mathbb{I}(s = R) + \tau_t \cdot \mathbb{I}(\text{ckpt} = t) + \Gamma_t \cdot (\mathbb{I}(s = R) \times \mathbb{I}(\text{ckpt} = t)) + \varepsilon_{i,s,t,k}
\end{equation}
where $\alpha_i$ absorbs intrinsic problem difficulty, $\beta$ captures baseline state disparity, $\tau_t$ captures clean proficiency growth, and $\Gamma_t$ identifies the excess recovery adaptation. This within-problem blocked contrast completely eliminates between-subject selection bias.

\paragraph{Statistical Inference and Clustering.}
Because observations are nested within problem instances ($K=16$ replications across $2 \times 2 = 4$ cells per problem), observations share within-problem correlation. We employ a \textbf{Problem-Blocked Bootstrap} with $B=10,000$ resamples, resampling entire problem clusters $\{p_i\}_{i=1}^N$ with replacement. Analytical standard errors are confirmed via Huber-White cluster-robust variance estimators clustered on problem IDs.

\section{Controlled Endpoint Experiment (Study A)}

\subsection{Confirmatory Results \texorpdfstring{($t=0$ vs.\ $t=256$)}{(t=0 vs t=256)}}
Our primary confirmatory trial evaluates $N=454$ problems across $29,056$ rollouts. As reported in Table~\ref{tab:endpoint}, transition success from base model ($t=0$) to terminal RL checkpoint ($t=256$) increases by $+32.05\text{ pp}$ under the recovery condition ($\mu_{R,0} = 0.3834 \to \mu_{R,256} = 0.7039$) and by $+20.29\text{ pp}$ under the control condition ($\mu_{C,0} = 0.3892 \to \mu_{C,256} = 0.5921$).

The resulting interaction contrast is:
\begin{equation}
\Gamma_{256} = +0.3205 - (+0.2029) = \mathbf{+0.1176} \quad (\mathbf{+11.76\text{ percentage points}})
\end{equation}
with a $95\%$ problem-blocked bootstrap confidence interval of $[+0.0955, +0.1400]$ ($z = 10.33, p < 10^{-24}$).

\input{tables/table1_endpoint.tex}

\begin{figure}[t]
\centering
\includegraphics[width=0.85\linewidth]{figures/figure3_cell_means.pdf}
\caption{\textbf{Primary factorial cell means ($2 \times 2$) in Study A.} The slope under Recovery ($+32.05\text{ pp}$) significantly exceeds the slope under Control ($+20.29\text{ pp}$), yielding $\Gamma_{256} = +11.76\text{ pp}$.}
\label{fig:cell_means}
\end{figure}

\subsection{Decoupling of Recovery Gains from Forward Proficiency}
To examine whether recovery adaptation is an incidental byproduct of general problem-solving growth, we analyze the problem-level gains $\Delta_{R,i}$ and $\Delta_{C,i}$ across the $454$ problem instances:
\begin{equation}
r(\Delta_{R,i}, \Delta_{C,i}) = \mathbf{0.0794}, \quad p = \mathbf{0.0911}, \quad R^2 = \mathbf{0.63\%}
\end{equation}
The correlation is near zero and statistically non-significant. A model checkpoint's ability to recover from intermediate fallacies constitutes an independent capability dimension that cannot be inferred from clean forward accuracy.

\subsection{Strict Decontamination Sensitivity Cohort \texorpdfstring{($N=388$)}{(N=388)}}
To verify that results are not driven by training set memorization, we filtered the benchmark against DeepScaleR's $40,000$-problem training set using 5-gram Jaccard matching ($J \ge 0.55$) and SequenceMatcher edit distance ($E \ge 0.70$). Removing 66 borderline/adversarial cases yields the strictly decontaminated cohort of $N_{\text{Strict}} = 388$ problems ($24,832$ rollouts).

Re-evaluating the estimand yields:
\begin{equation}
\Gamma_{256,\text{Strict}} = \mathbf{+0.1160} \quad (\mathbf{+11.60\text{ percentage points}})
\end{equation}
with $95\%$ bootstrap CI $[+0.0913, +0.1408]$ ($p < 10^{-20}$). The estimate changes by only $0.16\text{ pp}$ (a $1.3\%$ relative difference), confirming that benchmark contamination does not explain the recovery effect.

\subsection{Subgroup Stratification by Operator and Difficulty}
As shown in Table~\ref{tab:operators}, positive recovery interaction is observed across all error operators, with the largest contrast on arithmetic and constant corruptions ($\Gamma_{256} = +0.1432$). Furthermore, Table~\ref{tab:difficulty} indicates that recovery gains $\Delta_R$ remain steady ($\approx +31\text{ to }+39\text{ pp}$) across all difficulty tiers $L_1 \dots L_5$, refuting ceiling-effect artifacts.

\input{tables/table4_operators.tex}
\input{tables/table5_difficulty.tex}

\section{Longitudinal Checkpoint Trajectory \& Multiplicity Control}

\subsection{Empirical Nine-Checkpoint Trajectory}
We tracked $\Gamma_t$ across nine sequential post-training checkpoints $t \in \{0, 32, 64, 96, 128, 160, 192, 224, 256\}$. As reported in Table~\ref{tab:trajectory} and Figure~\ref{fig:trajectory}, the empirical trajectory vector is:
\begin{equation}
\mathbf{\Gamma} = [0.0000, +0.0333, +0.0337, +0.0774, +0.0748, +0.0598, +0.0976, +0.0950, +0.1176]
\end{equation}

\input{tables/table2_trajectory.tex}

\begin{figure}[t]
\centering
\includegraphics[width=0.88\linewidth]{figures/figure2_trajectory.pdf}
\caption{\textbf{Nine-checkpoint longitudinal trajectory of $\Gamma_t$.} Blue points denote unconstrained empirical estimates with $95\%$ confidence intervals. The dashed red line illustrates the fitted PAVA order-restricted monotonic envelope ($\bar{\chi}^2 = 42.18, p < 0.0001$).}
\label{fig:trajectory}
\end{figure}

\subsection{Non-Monotonicity and Order-Restricted Regression (PAVA)}
Examination of adjacent checkpoint steps reveals three local negative fluctuations in the unconstrained sample trajectory: Step $96 \to 128$ ($-0.26\text{ pp}$), Step $128 \to 160$ ($\mathbf{-1.50\text{ pp}}$), and Step $192 \to 224$ ($-0.26\text{ pp}$). Strict pairwise sample monotonicity across all adjacent checkpoints is therefore empirically refuted.

These fluctuations reflect intermediate Monte Carlo sampling variance resulting from resource-constrained replication allocations ($K=2/3$ rollouts at intermediate steps). Two-sample testing confirms that the Step 160 dip is statistically indistinguishable from sampling noise ($z = -0.595, p = 0.276$).

Applying the Pool Adjacent Violators Algorithm (PAVA)~\cite{Robertson1988Order} yields the least-squares monotonic envelope:
\begin{equation}
\mathbf{\Gamma}^* = [0.0000, +0.0333, +0.0337, +0.0707, +0.0707, +0.0707, +0.0963, +0.0963, +0.1176]
\end{equation}
A likelihood ratio test of the order-restricted isotonic model against the null hypothesis $\Gamma_t = 0$ rejects the null with $\bar{\chi}^2 = 42.18, p < 0.0001$. Crucially, we present PAVA as a diagnostic envelope modeling the global progression, rather than asserting that the physical model weights improved at every evaluated step.

\subsection{Multiplicity Adjustments \& Emergence Boundaries}
To prevent Type I error inflation across the eight post-baseline checkpoints, we apply Bonferroni, Holm-Bonferroni~\cite{Holm1979Simple}, and Benjamini-Hochberg FDR~\cite{Benjamini1995Controlling} adjustments:
\begin{itemize}
    \item \textbf{Step 32}: The contrast is nominally positive ($\Gamma_{32} = +0.0333, p_{\text{unadj}} = 0.0190$). However, under Bonferroni correction ($m=8$), $p_{\text{Bonf}} = 0.1520$ (and $p_{\text{Bonf}} = 0.0760$ under $m=4$), with simultaneous $95\%$ CI $[-0.0055, +0.0721]$ crossing zero. We transparently classify Step 32 as a \emph{nominally suggestive signal} rather than confirmed emergence.
    \item \textbf{Step 96}: By Step 96, state-dependent recovery achieves robust, family-wise confirmed statistical significance ($\Gamma_{96} = +0.0774, p_{\text{unadj}} = 1.71 \times 10^{-5}, p_{\text{Bonf}} = 0.0001 < 0.001$), with simultaneous $95\%$ CI $[+0.0282, +0.1266]$.
\end{itemize}

\section{Test-Time Compute \& Sequence-Length Invariance}

A critical potential confounder in reasoning evaluation is test-time search expansion: post-trained models might achieve higher recovery simply by generating longer reasoning traces.

\paragraph{Empirical Sequence Length Accounting.}
Across all $29,056$ rollouts, generated token lengths under the Recovery condition are virtually identical between baseline and terminal checkpoints:
\begin{itemize}
    \item Step 0 Recovery Tokens: $\text{Mean} = 482.4747$ tokens ($\text{SD} = 30.0882$, $\text{Median} = 490.0$).
    \item Step 256 Recovery Tokens: $\text{Mean} = 482.4953$ tokens ($\text{SD} = 30.4684$, $\text{Median} = 489.0$).
    \item Difference: $\Delta = \mathbf{+0.0206\text{ tokens}}$ ($t = -0.0411, p = 0.9672$).
\end{itemize}

\paragraph{Two One-Sided Tests (TOST) Equivalence Proof.}
A high $p$-value ($p = 0.9672$) demonstrates failure to detect a difference, not positive proof of equivalence. We therefore conduct Two One-Sided Tests (TOST) of statistical equivalence~\cite{Schuirmann1987Comparison}. In chain-of-thought generation, verbal reflection or re-derivation requires at least $15$ to $50$ additional tokens. Specifying a conservative equivalence margin of $\delta = \pm 1.0$ token:
\begin{equation}
t_1 = \frac{+0.0206 - (-1.0)}{0.5024} = +2.032, \quad t_2 = \frac{1.0 - 0.0206}{0.5024} = +1.950 \implies p_{\text{TOST}} = \mathbf{0.0256} < 0.05
\end{equation}
At margin $\delta = \pm 2.0$ tokens, $p_{\text{TOST}} = 4.10 \times 10^{-5}$. On paired problem-level means ($N=454$), TOST confirms equivalence at $\delta = \pm 1.0$ token with $p_{\text{TOST}} = \mathbf{0.0234}$. The recovery improvement occurs under strictly invariant test-time compute expenditure.

\section{Autonomous Natural Post-Error Recovery (Study B)}

To verify ecological validity beyond artificial prompt perturbations, Study B evaluates $3,200$ unprompted, unperturbed rollouts from Step 256 ($N=200$ problems, $K=16$). 

\input{tables/table3_natural.tex}

\begin{figure}[t]
\centering
\includegraphics[width=0.80\linewidth]{figures/figure4_funnel.pdf}
\caption{\textbf{Natural reasoning error and autonomous recovery funnel (Study B).} Out of $3,200$ unperturbed rollouts, $582$ contain endogenous errors ($\text{NEI}=18.19\%$), from which the model autonomously recovers in $180$ cases ($\text{NRR}=30.93\%$).}
\label{fig:funnel}
\end{figure}

As reported in Table~\ref{tab:natural} and Figure~\ref{fig:funnel}:
\begin{enumerate}
    \item \textbf{Natural Error Incidence ($\text{NEI}$)}: Intermediate reasoning deductions were audited by step verifiers, detecting endogenous reasoning errors in $582$ rollouts ($\text{NEI} = \mathbf{18.19\%}$, $95\%$ blocked CI $[16.84\%, 19.50\%]$).
    \item \textbf{Natural Post-Error Recovery Rate ($\text{NRR}$)}: Among the $582$ qualifying error episodes, the model autonomously redirected its reasoning trajectory and reached the verified terminal answer in $180$ cases ($\text{NRR} = \mathbf{30.93\%}$, $95\%$ blocked CI $[27.19\%, 34.82\%]$).
\end{enumerate}
This demonstrates that recovery is not an artifact of artificial prompt formatting; post-trained models autonomously self-correct in nearly one-third of natural failure episodes.

\section{Prospective Controlled Intervention Protocol (Study C)}

To isolate the precise causal mechanism of state-dependent recovery against alternative prompt paradigms, we designed a prospective within-subject $4 \times 100$ factorial experiment on the canonical \texttt{Qwen/Qwen2.5-Math-7B-Instruct} model (pinned commit \texttt{ef9926d}).

\subsection{Experimental Invariants and Cohort Design}
A cohort of $N=100$ challenging competition math problems from MATH-500 is evaluated under four orthogonal, strictly controlled prompt conditions ($400$ total frozen requests, SHA-256: \texttt{3a8caf3b26...}):
\begin{enumerate}
    \item \textbf{Condition A (Zero-Shot Baseline)}: Unassisted standard step-by-step reasoning.
    \item \textbf{Condition D (Forced Restart)}: Conditioned on a prefix containing an injected error at step $k$, followed by an instruction to discard all prior steps and restart from scratch.
    \item \textbf{Condition E (Explicit Repair)}: Conditioned on the identical prefix, followed by an instruction to identify and surgically repair the error.
    \item \textbf{Condition F (State-Preserving Recovery)}: Conditioned on the identical prefix, followed by an explicit directive to preserve all valid earlier steps verbatim and correct only the flawed assertion.
\end{enumerate}

\paragraph{Information Parity and Protocol Freeze.}
As verified through an exhaustive pre-experiment audit, all four conditions share identical problem stems, Conditions D, E, and F share identical prefix contexts up to the final instructional boundary, and zero ground-truth answers are leaked into prompts. Preliminary feasibility on local quantized models (1.5B) confirms the operational viability of the protocol. In accordance with strict pre-registration protocols, full physical GPU execution of the 400-request cohort on NVIDIA RTX 4090 hardware is subject to a two-stage preflight calibration gate; primary statistical estimands ($\tau_{F, D}$, $\tau_{F, E}$, and $\tau_{F, A}$) remain pending post-calibration execution.

\section{Discussion, Limitations \& Failure Horizon}

\paragraph{The Persistent 29.61\% Failure Regime.}
While RLVR dramatically improves recovery ($\mu_{R,0}=38.34\% \to \mu_{R,256}=70.39\%$), the model at Step 256 fails to recover from $29.61\%$ of invalid intermediate states ($1 - 0.7039 = 0.2961$). In single-pass autoregressive generation, certain arithmetic and algebraic fallacies function as catastrophic absorbing states: once generated, the conditioning history exerts strong probability mass on downstream inconsistent deductions. Overcoming this residual failure regime likely requires explicit multi-turn backtracking or external execution environments~\cite{Kumar2024Training}.

\paragraph{Threats to Validity \& Scope Boundaries.}
\begin{enumerate}
    \item \emph{Model Lineage}: Our empirical trajectory is evaluated on DeepScaleR-7B / Qwen-2.5-Math-7B. While highly powered ($29,056$ rollouts), cross-architecture generalizability to other model families (e.g., Llama-3, DeepSeek-V3) requires future empirical replication.
    \item \emph{Domain Specificity}: Evaluations are restricted to formal competition mathematics (MATH-500). Extending state-dependent evaluation to open-ended code generation and interactive tool use remains an important open direction.
\end{enumerate}

\section{Conclusion}
StateShift demonstrates that aggregate benchmark accuracy obscures critical state-dependent adaptation in post-trained language models. By isolating the State-Matched Interaction Contrast, we provide rigorous empirical proof that outcome-based reinforcement learning induces an autonomous, compute-invariant recovery capability ($\Gamma_{256} = +11.76\text{ pp}, p < 10^{-24}$) that is uncorrelated with clean forward accuracy gains. StateShift establishes an open diagnostic evaluation framework for auditing reasoning reliability in autonomous AI agents.

\bibliographystyle{plain}
\bibliography{references}

\end{document}
